% ============================================================
% MATERIAL APARCADO PARA EL TERCER PAPER (siembra/hibridación)
% Extraído de paper_gp/main.tex el 2026-07-24 (cirugía: el paper
% del GP queda autocontenido, sin comparación con el DRL).
%
% ATENCIÓN: los números aquí son los de la convención VIEJA del
% makespan (pre-fix). Antes de reutilizar: pasada con los valores
% corregidos (ver paper_gp/AUDIT_MAKESPAN_FIX.md):
%   - GP single-shot 18.5 -> 18.6 ; tuned 17.67 -> 17.7
%   - best-of-1024 desde pools corregidos: GP 13.69, DRL 13.08
%     (la brecha publicada 14.1 vs 12.7 = 1.4 pts pasa a 0.6 pts)
%   - DRL greedy 19.4: pendiente de auditar (rollout de checkpoint)
% ============================================================

% ---------- Contribución (era la 4 de la intro del paper GP) ----------
% \item A controlled comparison against a deep reinforcement learning
%   constructive policy trained on the same instances (companion
%   work~\cite{DiazRL2026}), with training and inference budgets matched:
%   single-shot, the interpretable rule matches the neural policy; under
%   equal sampling budgets the learned sampling distribution of the policy
%   retains an advantage (Section~\ref{sec:vs-rl}). Head-to-head
%   evaluations of the two learning paradigms under common protocols are
%   precisely what the recent survey of learned optimisation for the job
%   shop identifies as missing~\cite{Xu2025}; ours is, in addition, the
%   first in a scheduling setting with interval uncertainty.

% ---------- Contraste de tuning (era el cierre de §tuning) ----------
% This contrasts instructively with the neural policy of the companion
% work, for which operating-fidelity tuning did \emph{not} improve on the
% literature defaults at all: the flatter, lower-dimensional search space
% of the GP admits a small tuning gain where the co-adapted neural
% optimizer does not.

% ---------- Sección completa (era §vs-rl del paper GP) ----------
\subsection{Comparison with a neural constructive policy}
\label{sec:vs-rl}
GP rules and DRL policies are the two dominant learning paradigms for
constructive scheduling, yet controlled comparisons between them---same
instances, same evaluator, matched budgets---remain scarce, a gap
highlighted by the recent survey of Xu et al.~\cite{Xu2025}. This section
provides such a comparison for the IJSP. The companion DRL
policy~\cite{DiazRL2026} is trained on the same four
instances with a comparable wall-clock budget and evaluated with the same
protocol. Single-shot, the evolved rule \emph{matches} the neural policy
($18.5\%$ vs.\ $19.4\%$ greedy); at matched best-of-1024 the policy leads
($12.7\%$ vs.\ $14.1\%$), converting samples into quality faster
($-6.7$ vs.\ $-4.4$ points per ten doublings). The net advantage of deep
RL in this domain is therefore the quality of its learned sampling
distribution rather than a stronger point heuristic---while the GP rule
offers interpretability, zero-dependency deployment, and training two
orders of magnitude cheaper. With the tuned configuration of
Section~\ref{sec:tuning} the single-shot advantage of the rule widens
slightly ($17.7\%$ versus $19.4\%$). Table~\ref{tab:matched} summarizes
the matched-budget comparison.

\begin{table}[t]
\centering\small
\caption{Matched-budget comparison over the 70 instances: mean RE (\%).
Training budgets are comparable (GP: ${\sim}20{,}400$ rollouts,
${\approx}49$ min; DRL: $4{,}000$ episodes, ${\approx}80$ min; identical
training instances).}
\label{tab:matched}
\begin{tabular}{lcc}
\toprule
Method & best-of-1 & best-of-1024 \\
\midrule
GP rule (default) / GP-$\varepsilon$ & 18.5 & 14.1 \\
GP rule (tuned) & \textbf{17.7} & --- \\
DRL policy (greedy / sampled) & 19.4 & \textbf{12.7} \\
\bottomrule
\end{tabular}
\end{table}

% ---------- Conclusión (era la (iv) del paper GP) ----------
% (iv)~Against that neural policy under matched budgets, the interpretable
% rule matches or beats it single-shot, and cedes ground only when both
% methods are granted large sampling budgets, where the learned sampling
% distribution of the policy is more efficient than uniform randomization.

% ---------- Bibliografía asociada ----------
% \bibitem{DiazRL2026}
% H.~D\'iaz et al.
% \newblock Size-invariant deep reinforcement learning for the interval job
%   shop scheduling problem.
% \newblock (estado por decidir al preparar el tercer paper)
